% USE PDFLATEX WITH THIS DOCUMENT. If you wish to use ordinary latex (why?) then convert the figures in figures/ to eps.
% If you are using Chalmers linux system, you might need to change to a newer tex dist.
% Do so with: vcs-what latex 
% look up the newest version and select it, a.t.m that would be: vcs-select -p texlive-20080816
% I recommend using kile in linux. Make it available by adding the unsup software 
% echo unsup ~/.vcs4/pathsetup
\documentclass[a4paper,12pt,twoside]{article}
\usepackage[english]{babel}
%\usepackage[latin1]{inputenc} % Use the same as the encoding of the textfile
\usepackage[utf8]{inputenc} % But this is probably the best choice. This doument is currently in utf8 (so is all other files)
\usepackage{amsmath,amsfonts,amsthm,mathtools} % Math. Look up the amsmath manual (google) for many great examples.
\usepackage{graphicx,placeins,float} % Figures and placements, float must be loaded before hypperref
\usepackage[colorlinks,citecolor=blue]{hyperref} % Warning this package conflicts with the package algorithm
\usepackage{fancyhdr,url,tikz} % Other usefull stuff. tikz is good for making drawings (but difficult)
\usepackage{listings} % Source code
\usepackage{a4wide} % Gives a wider page
\usepackage{mycommands} % Here you can define your personal favourites. See mycommands.cls
\usepackage{pgfplots} % Another nice library for plots

% mysubtitle, cover* and divisionnameB may be uncommented if nonexistant.
\newcommand{\mytitle}{Stability of High Speed Train under Aero\-dynamic Excitations}
%\newcommand{\mysubtitle}{Some subtitle}
\newcommand{\writtenmonth}{April}
\newcommand{\writtenyear}{2009}
\newcommand{\authors}{Erik Bjerklund\\Mikael \"Ohman}
\newcommand{\authorsc}{Erik Bjerklund, Mikael \"Ohman} % again with comma
\newcommand{\YYYYNN}{2009:03}
\newcommand{\ISSN}{1652-8557}
\newcommand{\covercaption}{Train leaving a tunnel and is hit by a strong side wind.}
\newcommand{\coverfigure}{figures/COVER_93_iso.png}
\newcommand{\departmentname}{Applied Mechanics}
\newcommand{\divisionnameA}{Dynamics}
\newcommand{\divisionnameB}{Fluid Mechanics}
\newcommand{\nameofprinters}{Chalmers Reproservice}
\newcommand{\nameofprogramme}{Solid and Fluid Mechanics}
\newcommand{\keywords}{Moving mesh, Wind loads, Coupling, Vibration dynamics, CFD, High speed train, Comfort, Stability}
% This is just to keep the swedish letters at one place in case there are problems with encoding
\newcommand{\city}{G\"oteborg} 
\newcommand{\country}{Sweden} 
\newcommand{\university}{Chalmers University of Technology}
\newcommand{\thesis}{Master's Thesis}
\newcommand{\telephone}{+ 46 (0)31-772 1000}
\newcommand{\postcode}{SE-412 96}
% End of input.

% Use section number first in numbering
\numberwithin{equation}{section}
\numberwithin{figure}{section}
\numberwithin{table}{section}

% Setting up the marginals a bit larger
\setlength{\oddsidemargin}{0mm}
\setlength{\textheight}{730pt} % default 609pt, ~710 if you dont use a4wide
\setlength{\headsep}{0pt} % default 25pt, since we have no headers in this document
\setlength{\headheight}{0pt} % 12 pt
\setlength{\voffset}{-0.4in} % default 0

% Set up headers to match styleguide
\pagestyle{fancy}
\renewcommand{\headrulewidth}{0pt}
\fancyhead{}
\fancyfoot{}
\fancyfoot[C]{\footnotesize \includegraphics[height=2.5mm]{figures/CHALMERS_black.pdf}, \textit{\departmentname}, \thesis\ \YYYYNN}
\fancyfoot[RO,LE]{\thepage}

% To use code use \lstinputlisting[langauge=matlab,style=mystyle]{somematlabfile.m}
% Here is a suitable style for code.
\lstdefinestyle{mystyle}{showstringspaces=false, basicstyle=\scriptsize\ttfamily,
        frame=shadowbox, breaklines=true, numbers=left, commentstyle=\color{gray},
        keywordstyle=\color{blue}\textbf, stringstyle=\color{red}}

% Metadata ion the PDF file (makes it searchable)
\hypersetup{pdfauthor={\authorsc},pdftitle={\mytitle},pdfsubject={\thesis},pdfkeywords={\keywords}}

% The following part is automatically generated, go to document start
\begin{document}
\thispagestyle{empty}

\newlength{\headerheight};
\newlength{\headeroffset};
\setlength{\headerheight}{0.215\paperheight}%~{6.4cm};
\setlength{\headeroffset}{0.125\headerheight};

%\begin{tikzpicture}[remember picture,overlay]

%  \node[yshift=-\headerheight] at (current page.north west)
%    {\begin{tikzpicture}[remember picture, overlay]
%      \clip[] (0,\headerheight) -- +(\paperwidth,0) -- (\paperwidth,0) -- (0.725\paperwidth,0)
%                .. controls (0.703\paperwidth,0) and (0.703\paperwidth,\headeroffset) ..
%                (0.68\paperwidth,\headeroffset) -- (0,\headeroffset) -- cycle;
%      \fill[black!15] (0,0) rectangle (\paperwidth,\headerheight);
%      \definecolor{logo_color}{gray}{1}
%      \node[anchor=west,xshift=0.1\paperwidth,yshift=0.511\headerheight]{\resizebox{0.63\paperwidth}{!}{\begin{pgfpicture}

\begin{pgfscope} % This is 11.994 cm in width, 1.341cm in height

\pgfsetstrokecolor{logo_color}
\pgfsetfillcolor{logo_color}

% C
\pgfpathmoveto{\pgfqpoint{1.481cm}{0.488cm}}
\pgfpathcurveto{\pgfqpoint{1.385cm}{0.138cm}}{\pgfqpoint{1.046cm}{0.027cm}}{\pgfqpoint{0.742cm}{0.027cm}}
\pgfpathcurveto{\pgfqpoint{0.304cm}{0.027cm}}{\pgfqpoint{0cm}{0.256cm}}{\pgfqpoint{0cm}{0.687cm}}
\pgfpathcurveto{\pgfqpoint{0cm}{1.152cm}}{\pgfqpoint{0.37cm}{1.34cm}}{\pgfqpoint{0.76cm}{1.34cm}}
\pgfpathcurveto{\pgfqpoint{0.901cm}{1.34cm}}{\pgfqpoint{1.273cm}{1.307cm}}{\pgfqpoint{1.437cm}{1.011cm}}
\pgfpathlineto{\pgfqpoint{1.128cm}{0.914cm}}
\pgfpathcurveto{\pgfqpoint{1.095cm}{0.966cm}}{\pgfqpoint{0.997cm}{1.102cm}}{\pgfqpoint{0.758cm}{1.102cm}}
\pgfpathcurveto{\pgfqpoint{0.478cm}{1.102cm}}{\pgfqpoint{0.321cm}{0.919cm}}{\pgfqpoint{0.321cm}{0.683cm}}
\pgfpathcurveto{\pgfqpoint{0.321cm}{0.414cm}}{\pgfqpoint{0.535cm}{0.271cm}}{\pgfqpoint{0.769cm}{0.271cm}}
\pgfpathcurveto{\pgfqpoint{1.036cm}{0.271cm}}{\pgfqpoint{1.116cm}{0.411cm}}{\pgfqpoint{1.156cm}{0.486cm}}
\pgfpathlineto{\pgfqpoint{1.481cm}{0.488cm}}
\pgfusepath{fill}

% H
\pgfpathmoveto{\pgfqpoint{1.616cm}{1.302cm}}
\pgfpathlineto{\pgfqpoint{1.925cm}{1.302cm}}
\pgfpathlineto{\pgfqpoint{1.925cm}{0.827cm}}
\pgfpathlineto{\pgfqpoint{2.656cm}{0.827cm}}
\pgfpathlineto{\pgfqpoint{2.656cm}{1.302cm}}
\pgfpathlineto{\pgfqpoint{2.965cm}{1.302cm}}
\pgfpathlineto{\pgfqpoint{2.965cm}{0.069cm}}
\pgfpathlineto{\pgfqpoint{2.656cm}{0.069cm}}
\pgfpathlineto{\pgfqpoint{2.656cm}{0.598cm}}
\pgfpathlineto{\pgfqpoint{1.925cm}{0.598cm}}
\pgfpathlineto{\pgfqpoint{1.925cm}{0.069cm}}
\pgfpathlineto{\pgfqpoint{1.616cm}{0.069cm}}
\pgfpathlineto{\pgfqpoint{1.616cm}{1.302cm}}
\pgfusepath{fill} 

%A
\pgfpathmoveto{\pgfqpoint{3.707cm}{1.302cm}}
\pgfpathlineto{\pgfqpoint{4.041cm}{1.302cm}}
\pgfpathlineto{\pgfqpoint{4.637cm}{0.069cm}}
\pgfpathlineto{\pgfqpoint{4.291cm}{0.069cm}}
\pgfpathlineto{\pgfqpoint{4.168cm}{0.338cm}}
\pgfpathlineto{\pgfqpoint{3.568cm}{0.338cm}}
\pgfpathlineto{\pgfqpoint{3.449cm}{0.069cm}}
\pgfpathlineto{\pgfqpoint{3.112cm}{0.069cm}}
\pgfpathlineto{\pgfqpoint{3.707cm}{1.302cm}}
\pgfpathclose
\pgfpathmoveto{\pgfqpoint{3.903cm}{0.923cm}}
\pgfpathcurveto{\pgfqpoint{3.889cm}{0.955cm}}{\pgfqpoint{3.876cm}{0.985cm}}{\pgfqpoint{3.871cm}{1.021cm}}
\pgfpathcurveto{\pgfqpoint{3.866cm}{0.986cm}}{\pgfqpoint{3.851cm}{0.954cm}}{\pgfqpoint{3.837cm}{0.922cm}}
\pgfpathlineto{\pgfqpoint{3.668cm}{0.556cm}}
\pgfpathlineto{\pgfqpoint{4.066cm}{0.556cm}}
\pgfpathlineto{\pgfqpoint{3.903cm}{0.923cm}}
\pgfusepath{fill}

% L
\pgfpathmoveto{\pgfqpoint{4.796cm}{1.302cm}}
\pgfpathlineto{\pgfqpoint{5.105cm}{1.302cm}}
\pgfpathlineto{\pgfqpoint{5.105cm}{0.296cm}}
\pgfpathlineto{\pgfqpoint{5.728cm}{0.296cm}}
\pgfpathlineto{\pgfqpoint{5.728cm}{0.069cm}}
\pgfpathlineto{\pgfqpoint{4.796cm}{0.069cm}}
\pgfpathlineto{\pgfqpoint{4.796cm}{1.302cm}}
\pgfusepath{fill}

% M
\pgfpathmoveto{\pgfqpoint{5.914cm}{1.302cm}}
\pgfpathlineto{\pgfqpoint{6.359cm}{1.302cm}}
\pgfpathlineto{\pgfqpoint{6.687cm}{0.532cm}}
\pgfpathcurveto{\pgfqpoint{6.721cm}{0.451cm}}{\pgfqpoint{6.73cm}{0.418cm}}{\pgfqpoint{6.743cm}{0.364cm}}
\pgfpathcurveto{\pgfqpoint{6.756cm}{0.418cm}}{\pgfqpoint{6.767cm}{0.451cm}}{\pgfqpoint{6.8cm}{0.532cm}}
\pgfpathlineto{\pgfqpoint{7.122cm}{1.302cm}}
\pgfpathlineto{\pgfqpoint{7.564cm}{1.302cm}}
\pgfpathlineto{\pgfqpoint{7.564cm}{0.069cm}}
\pgfpathlineto{\pgfqpoint{7.272cm}{0.069cm}}
\pgfpathlineto{\pgfqpoint{7.272cm}{0.937cm}}
\pgfpathcurveto{\pgfqpoint{7.272cm}{0.984cm}}{\pgfqpoint{7.279cm}{1.01cm}}{\pgfqpoint{7.285cm}{1.047cm}}
\pgfpathcurveto{\pgfqpoint{7.272cm}{1.01cm}}{\pgfqpoint{7.263cm}{0.984cm}}{\pgfqpoint{7.243cm}{0.937cm}}
\pgfpathlineto{\pgfqpoint{6.881cm}{0.069cm}}
\pgfpathlineto{\pgfqpoint{6.6cm}{0.069cm}}
\pgfpathlineto{\pgfqpoint{6.231cm}{0.937cm}}
\pgfpathcurveto{\pgfqpoint{6.211cm}{0.984cm}}{\pgfqpoint{6.205cm}{1.01cm}}{\pgfqpoint{6.195cm}{1.047cm}}
\pgfpathcurveto{\pgfqpoint{6.198cm}{1.01cm}}{\pgfqpoint{6.205cm}{0.984cm}}{\pgfqpoint{6.205cm}{0.937cm}}
\pgfpathlineto{\pgfqpoint{6.205cm}{0.069cm}}
\pgfpathlineto{\pgfqpoint{5.914cm}{0.069cm}}
\pgfpathlineto{\pgfqpoint{5.914cm}{1.302cm}}
\pgfusepath{fill}

% E
\pgfpathmoveto{\pgfqpoint{7.85cm}{1.302cm}}
\pgfpathlineto{\pgfqpoint{8.999cm}{1.302cm}}
\pgfpathlineto{\pgfqpoint{8.999cm}{1.08cm}}
\pgfpathlineto{\pgfqpoint{8.159cm}{1.08cm}}
\pgfpathlineto{\pgfqpoint{8.159cm}{0.802cm}}
\pgfpathlineto{\pgfqpoint{8.899cm}{0.802cm}}
\pgfpathlineto{\pgfqpoint{8.899cm}{0.584cm}}
\pgfpathlineto{\pgfqpoint{8.159cm}{0.584cm}}
\pgfpathlineto{\pgfqpoint{8.159cm}{0.289cm}}
\pgfpathlineto{\pgfqpoint{9.034cm}{0.289cm}}
\pgfpathlineto{\pgfqpoint{9.034cm}{0.069cm}}
\pgfpathlineto{\pgfqpoint{7.85cm}{0.069cm}}
\pgfpathlineto{\pgfqpoint{7.85cm}{1.302cm}}
\pgfusepath{fill}

% R
\pgfpathmoveto{\pgfqpoint{9.266cm}{1.301cm}}
\pgfpathlineto{\pgfqpoint{10.128cm}{1.301cm}}
\pgfpathcurveto{\pgfqpoint{10.219cm}{1.301cm}}{\pgfqpoint{10.566cm}{1.29cm}}{\pgfqpoint{10.566cm}{0.933cm}}
\pgfpathcurveto{\pgfqpoint{10.566cm}{0.72cm}}{\pgfqpoint{10.439cm}{0.612cm}}{\pgfqpoint{10.246cm}{0.583cm}}
\pgfpathlineto{\pgfqpoint{10.62cm}{0.068cm}}
\pgfpathlineto{\pgfqpoint{10.265cm}{0.068cm}}
\pgfpathlineto{\pgfqpoint{9.927cm}{0.564cm}}
\pgfpathlineto{\pgfqpoint{9.575cm}{0.564cm}}
\pgfpathlineto{\pgfqpoint{9.575cm}{0.068cm}}
\pgfpathlineto{\pgfqpoint{9.266cm}{0.068cm}}
\pgfpathlineto{\pgfqpoint{9.266cm}{1.301cm}}
\pgfpathclose
\pgfpathmoveto{\pgfqpoint{9.576cm}{0.782cm}}
\pgfpathlineto{\pgfqpoint{10.068cm}{0.782cm}}
\pgfpathcurveto{\pgfqpoint{10.134cm}{0.782cm}}{\pgfqpoint{10.253cm}{0.791cm}}{\pgfqpoint{10.253cm}{0.929cm}}
\pgfpathcurveto{\pgfqpoint{10.253cm}{1.064cm}}{\pgfqpoint{10.164cm}{1.085cm}}{\pgfqpoint{10.054cm}{1.085cm}}
\pgfpathlineto{\pgfqpoint{9.576cm}{1.085cm}}
\pgfpathlineto{\pgfqpoint{9.576cm}{0.782cm}}
\pgfusepath{fill}


% S
\pgfpathmoveto{\pgfqpoint{10.956cm}{0.445cm}}
\pgfpathcurveto{\pgfqpoint{11.028cm}{0.297cm}}{\pgfqpoint{11.205cm}{0.255cm}}{\pgfqpoint{11.373cm}{0.258cm}}
\pgfpathcurveto{\pgfqpoint{11.448cm}{0.259cm}}{\pgfqpoint{11.684cm}{0.276cm}}{\pgfqpoint{11.681cm}{0.426cm}}
\pgfpathcurveto{\pgfqpoint{11.68cm}{0.515cm}}{\pgfqpoint{11.593cm}{0.534cm}}{\pgfqpoint{11.527cm}{0.545cm}}
\pgfpathcurveto{\pgfqpoint{11.476cm}{0.554cm}}{\pgfqpoint{11.222cm}{0.589cm}}{\pgfqpoint{11.164cm}{0.599cm}}
\pgfpathcurveto{\pgfqpoint{11.051cm}{0.617cm}}{\pgfqpoint{10.758cm}{0.667cm}}{\pgfqpoint{10.758cm}{0.946cm}}
\pgfpathcurveto{\pgfqpoint{10.758cm}{1.249cm}}{\pgfqpoint{11.08cm}{1.34cm}}{\pgfqpoint{11.326cm}{1.34cm}}
\pgfpathcurveto{\pgfqpoint{11.612cm}{1.34cm}}{\pgfqpoint{11.831cm}{1.285cm}}{\pgfqpoint{11.981cm}{1.056cm}}
\pgfpathlineto{\pgfqpoint{11.689cm}{0.978cm}}
\pgfpathcurveto{\pgfqpoint{11.62cm}{1.103cm}}{\pgfqpoint{11.462cm}{1.121cm}}{\pgfqpoint{11.33cm}{1.12cm}}
\pgfpathcurveto{\pgfqpoint{11.194cm}{1.12cm}}{\pgfqpoint{11.055cm}{1.089cm}}{\pgfqpoint{11.057cm}{0.972cm}}
\pgfpathcurveto{\pgfqpoint{11.058cm}{0.892cm}}{\pgfqpoint{11.13cm}{0.872cm}}{\pgfqpoint{11.223cm}{0.853cm}}
\pgfpathcurveto{\pgfqpoint{11.274cm}{0.843cm}}{\pgfqpoint{11.526cm}{0.805cm}}{\pgfqpoint{11.584cm}{0.795cm}}
\pgfpathcurveto{\pgfqpoint{11.727cm}{0.77cm}}{\pgfqpoint{11.971cm}{0.723cm}}{\pgfqpoint{11.976cm}{0.452cm}}
\pgfpathcurveto{\pgfqpoint{11.983cm}{0.056cm}}{\pgfqpoint{11.539cm}{0.033cm}}{\pgfqpoint{11.358cm}{0.031cm}}
\pgfpathcurveto{\pgfqpoint{11.156cm}{0.028cm}}{\pgfqpoint{10.84cm}{0.057cm}}{\pgfqpoint{10.676cm}{0.366cm}}
\pgfpathlineto{\pgfqpoint{10.956cm}{0.445cm}}
\pgfusepath{fill}

\end{pgfscope}
\end{pgfpicture}
}}; % 13.2 cm width
%      \node[anchor=west,xshift=0.745\paperwidth,yshift=0.516\headerheight]{\includegraphics[width=0.171\paperwidth]{figures/Avancez.pdf}}; % 3.6cm
%      \end{tikzpicture}};
% \end{tikzpicture}

\begin{tikzpicture}[remember picture,overlay]
 \node[yshift=-6.4cm] at (current page.north west)
   {\begin{tikzpicture}[remember picture, overlay]
     % First draw the grid and then the Logo and avancez logo.
     \clip[draw] (0,\headerheight) -- +(\paperwidth,0) -- (\paperwidth,0) -- (0.725\paperwidth,0)
                .. controls (0.703\paperwidth,0) and (0.703\paperwidth,\headeroffset) ..
                (0.68\paperwidth,\headeroffset) -- (0,\headeroffset) -- cycle;
     \draw[step=5mm,black] (0,0) grid (\paperwidth,6.4cm);
     \definecolor{logo_color}{gray}{0}
     \node[anchor=west,xshift=0.1\paperwidth,yshift=0.511\headerheight]{\resizebox{0.63\paperwidth}{!}{\begin{pgfpicture}

\begin{pgfscope} % This is 11.994 cm in width, 1.341cm in height

\pgfsetstrokecolor{logo_color}
\pgfsetfillcolor{logo_color}

% C
\pgfpathmoveto{\pgfqpoint{1.481cm}{0.488cm}}
\pgfpathcurveto{\pgfqpoint{1.385cm}{0.138cm}}{\pgfqpoint{1.046cm}{0.027cm}}{\pgfqpoint{0.742cm}{0.027cm}}
\pgfpathcurveto{\pgfqpoint{0.304cm}{0.027cm}}{\pgfqpoint{0cm}{0.256cm}}{\pgfqpoint{0cm}{0.687cm}}
\pgfpathcurveto{\pgfqpoint{0cm}{1.152cm}}{\pgfqpoint{0.37cm}{1.34cm}}{\pgfqpoint{0.76cm}{1.34cm}}
\pgfpathcurveto{\pgfqpoint{0.901cm}{1.34cm}}{\pgfqpoint{1.273cm}{1.307cm}}{\pgfqpoint{1.437cm}{1.011cm}}
\pgfpathlineto{\pgfqpoint{1.128cm}{0.914cm}}
\pgfpathcurveto{\pgfqpoint{1.095cm}{0.966cm}}{\pgfqpoint{0.997cm}{1.102cm}}{\pgfqpoint{0.758cm}{1.102cm}}
\pgfpathcurveto{\pgfqpoint{0.478cm}{1.102cm}}{\pgfqpoint{0.321cm}{0.919cm}}{\pgfqpoint{0.321cm}{0.683cm}}
\pgfpathcurveto{\pgfqpoint{0.321cm}{0.414cm}}{\pgfqpoint{0.535cm}{0.271cm}}{\pgfqpoint{0.769cm}{0.271cm}}
\pgfpathcurveto{\pgfqpoint{1.036cm}{0.271cm}}{\pgfqpoint{1.116cm}{0.411cm}}{\pgfqpoint{1.156cm}{0.486cm}}
\pgfpathlineto{\pgfqpoint{1.481cm}{0.488cm}}
\pgfusepath{fill}

% H
\pgfpathmoveto{\pgfqpoint{1.616cm}{1.302cm}}
\pgfpathlineto{\pgfqpoint{1.925cm}{1.302cm}}
\pgfpathlineto{\pgfqpoint{1.925cm}{0.827cm}}
\pgfpathlineto{\pgfqpoint{2.656cm}{0.827cm}}
\pgfpathlineto{\pgfqpoint{2.656cm}{1.302cm}}
\pgfpathlineto{\pgfqpoint{2.965cm}{1.302cm}}
\pgfpathlineto{\pgfqpoint{2.965cm}{0.069cm}}
\pgfpathlineto{\pgfqpoint{2.656cm}{0.069cm}}
\pgfpathlineto{\pgfqpoint{2.656cm}{0.598cm}}
\pgfpathlineto{\pgfqpoint{1.925cm}{0.598cm}}
\pgfpathlineto{\pgfqpoint{1.925cm}{0.069cm}}
\pgfpathlineto{\pgfqpoint{1.616cm}{0.069cm}}
\pgfpathlineto{\pgfqpoint{1.616cm}{1.302cm}}
\pgfusepath{fill} 

%A
\pgfpathmoveto{\pgfqpoint{3.707cm}{1.302cm}}
\pgfpathlineto{\pgfqpoint{4.041cm}{1.302cm}}
\pgfpathlineto{\pgfqpoint{4.637cm}{0.069cm}}
\pgfpathlineto{\pgfqpoint{4.291cm}{0.069cm}}
\pgfpathlineto{\pgfqpoint{4.168cm}{0.338cm}}
\pgfpathlineto{\pgfqpoint{3.568cm}{0.338cm}}
\pgfpathlineto{\pgfqpoint{3.449cm}{0.069cm}}
\pgfpathlineto{\pgfqpoint{3.112cm}{0.069cm}}
\pgfpathlineto{\pgfqpoint{3.707cm}{1.302cm}}
\pgfpathclose
\pgfpathmoveto{\pgfqpoint{3.903cm}{0.923cm}}
\pgfpathcurveto{\pgfqpoint{3.889cm}{0.955cm}}{\pgfqpoint{3.876cm}{0.985cm}}{\pgfqpoint{3.871cm}{1.021cm}}
\pgfpathcurveto{\pgfqpoint{3.866cm}{0.986cm}}{\pgfqpoint{3.851cm}{0.954cm}}{\pgfqpoint{3.837cm}{0.922cm}}
\pgfpathlineto{\pgfqpoint{3.668cm}{0.556cm}}
\pgfpathlineto{\pgfqpoint{4.066cm}{0.556cm}}
\pgfpathlineto{\pgfqpoint{3.903cm}{0.923cm}}
\pgfusepath{fill}

% L
\pgfpathmoveto{\pgfqpoint{4.796cm}{1.302cm}}
\pgfpathlineto{\pgfqpoint{5.105cm}{1.302cm}}
\pgfpathlineto{\pgfqpoint{5.105cm}{0.296cm}}
\pgfpathlineto{\pgfqpoint{5.728cm}{0.296cm}}
\pgfpathlineto{\pgfqpoint{5.728cm}{0.069cm}}
\pgfpathlineto{\pgfqpoint{4.796cm}{0.069cm}}
\pgfpathlineto{\pgfqpoint{4.796cm}{1.302cm}}
\pgfusepath{fill}

% M
\pgfpathmoveto{\pgfqpoint{5.914cm}{1.302cm}}
\pgfpathlineto{\pgfqpoint{6.359cm}{1.302cm}}
\pgfpathlineto{\pgfqpoint{6.687cm}{0.532cm}}
\pgfpathcurveto{\pgfqpoint{6.721cm}{0.451cm}}{\pgfqpoint{6.73cm}{0.418cm}}{\pgfqpoint{6.743cm}{0.364cm}}
\pgfpathcurveto{\pgfqpoint{6.756cm}{0.418cm}}{\pgfqpoint{6.767cm}{0.451cm}}{\pgfqpoint{6.8cm}{0.532cm}}
\pgfpathlineto{\pgfqpoint{7.122cm}{1.302cm}}
\pgfpathlineto{\pgfqpoint{7.564cm}{1.302cm}}
\pgfpathlineto{\pgfqpoint{7.564cm}{0.069cm}}
\pgfpathlineto{\pgfqpoint{7.272cm}{0.069cm}}
\pgfpathlineto{\pgfqpoint{7.272cm}{0.937cm}}
\pgfpathcurveto{\pgfqpoint{7.272cm}{0.984cm}}{\pgfqpoint{7.279cm}{1.01cm}}{\pgfqpoint{7.285cm}{1.047cm}}
\pgfpathcurveto{\pgfqpoint{7.272cm}{1.01cm}}{\pgfqpoint{7.263cm}{0.984cm}}{\pgfqpoint{7.243cm}{0.937cm}}
\pgfpathlineto{\pgfqpoint{6.881cm}{0.069cm}}
\pgfpathlineto{\pgfqpoint{6.6cm}{0.069cm}}
\pgfpathlineto{\pgfqpoint{6.231cm}{0.937cm}}
\pgfpathcurveto{\pgfqpoint{6.211cm}{0.984cm}}{\pgfqpoint{6.205cm}{1.01cm}}{\pgfqpoint{6.195cm}{1.047cm}}
\pgfpathcurveto{\pgfqpoint{6.198cm}{1.01cm}}{\pgfqpoint{6.205cm}{0.984cm}}{\pgfqpoint{6.205cm}{0.937cm}}
\pgfpathlineto{\pgfqpoint{6.205cm}{0.069cm}}
\pgfpathlineto{\pgfqpoint{5.914cm}{0.069cm}}
\pgfpathlineto{\pgfqpoint{5.914cm}{1.302cm}}
\pgfusepath{fill}

% E
\pgfpathmoveto{\pgfqpoint{7.85cm}{1.302cm}}
\pgfpathlineto{\pgfqpoint{8.999cm}{1.302cm}}
\pgfpathlineto{\pgfqpoint{8.999cm}{1.08cm}}
\pgfpathlineto{\pgfqpoint{8.159cm}{1.08cm}}
\pgfpathlineto{\pgfqpoint{8.159cm}{0.802cm}}
\pgfpathlineto{\pgfqpoint{8.899cm}{0.802cm}}
\pgfpathlineto{\pgfqpoint{8.899cm}{0.584cm}}
\pgfpathlineto{\pgfqpoint{8.159cm}{0.584cm}}
\pgfpathlineto{\pgfqpoint{8.159cm}{0.289cm}}
\pgfpathlineto{\pgfqpoint{9.034cm}{0.289cm}}
\pgfpathlineto{\pgfqpoint{9.034cm}{0.069cm}}
\pgfpathlineto{\pgfqpoint{7.85cm}{0.069cm}}
\pgfpathlineto{\pgfqpoint{7.85cm}{1.302cm}}
\pgfusepath{fill}

% R
\pgfpathmoveto{\pgfqpoint{9.266cm}{1.301cm}}
\pgfpathlineto{\pgfqpoint{10.128cm}{1.301cm}}
\pgfpathcurveto{\pgfqpoint{10.219cm}{1.301cm}}{\pgfqpoint{10.566cm}{1.29cm}}{\pgfqpoint{10.566cm}{0.933cm}}
\pgfpathcurveto{\pgfqpoint{10.566cm}{0.72cm}}{\pgfqpoint{10.439cm}{0.612cm}}{\pgfqpoint{10.246cm}{0.583cm}}
\pgfpathlineto{\pgfqpoint{10.62cm}{0.068cm}}
\pgfpathlineto{\pgfqpoint{10.265cm}{0.068cm}}
\pgfpathlineto{\pgfqpoint{9.927cm}{0.564cm}}
\pgfpathlineto{\pgfqpoint{9.575cm}{0.564cm}}
\pgfpathlineto{\pgfqpoint{9.575cm}{0.068cm}}
\pgfpathlineto{\pgfqpoint{9.266cm}{0.068cm}}
\pgfpathlineto{\pgfqpoint{9.266cm}{1.301cm}}
\pgfpathclose
\pgfpathmoveto{\pgfqpoint{9.576cm}{0.782cm}}
\pgfpathlineto{\pgfqpoint{10.068cm}{0.782cm}}
\pgfpathcurveto{\pgfqpoint{10.134cm}{0.782cm}}{\pgfqpoint{10.253cm}{0.791cm}}{\pgfqpoint{10.253cm}{0.929cm}}
\pgfpathcurveto{\pgfqpoint{10.253cm}{1.064cm}}{\pgfqpoint{10.164cm}{1.085cm}}{\pgfqpoint{10.054cm}{1.085cm}}
\pgfpathlineto{\pgfqpoint{9.576cm}{1.085cm}}
\pgfpathlineto{\pgfqpoint{9.576cm}{0.782cm}}
\pgfusepath{fill}


% S
\pgfpathmoveto{\pgfqpoint{10.956cm}{0.445cm}}
\pgfpathcurveto{\pgfqpoint{11.028cm}{0.297cm}}{\pgfqpoint{11.205cm}{0.255cm}}{\pgfqpoint{11.373cm}{0.258cm}}
\pgfpathcurveto{\pgfqpoint{11.448cm}{0.259cm}}{\pgfqpoint{11.684cm}{0.276cm}}{\pgfqpoint{11.681cm}{0.426cm}}
\pgfpathcurveto{\pgfqpoint{11.68cm}{0.515cm}}{\pgfqpoint{11.593cm}{0.534cm}}{\pgfqpoint{11.527cm}{0.545cm}}
\pgfpathcurveto{\pgfqpoint{11.476cm}{0.554cm}}{\pgfqpoint{11.222cm}{0.589cm}}{\pgfqpoint{11.164cm}{0.599cm}}
\pgfpathcurveto{\pgfqpoint{11.051cm}{0.617cm}}{\pgfqpoint{10.758cm}{0.667cm}}{\pgfqpoint{10.758cm}{0.946cm}}
\pgfpathcurveto{\pgfqpoint{10.758cm}{1.249cm}}{\pgfqpoint{11.08cm}{1.34cm}}{\pgfqpoint{11.326cm}{1.34cm}}
\pgfpathcurveto{\pgfqpoint{11.612cm}{1.34cm}}{\pgfqpoint{11.831cm}{1.285cm}}{\pgfqpoint{11.981cm}{1.056cm}}
\pgfpathlineto{\pgfqpoint{11.689cm}{0.978cm}}
\pgfpathcurveto{\pgfqpoint{11.62cm}{1.103cm}}{\pgfqpoint{11.462cm}{1.121cm}}{\pgfqpoint{11.33cm}{1.12cm}}
\pgfpathcurveto{\pgfqpoint{11.194cm}{1.12cm}}{\pgfqpoint{11.055cm}{1.089cm}}{\pgfqpoint{11.057cm}{0.972cm}}
\pgfpathcurveto{\pgfqpoint{11.058cm}{0.892cm}}{\pgfqpoint{11.13cm}{0.872cm}}{\pgfqpoint{11.223cm}{0.853cm}}
\pgfpathcurveto{\pgfqpoint{11.274cm}{0.843cm}}{\pgfqpoint{11.526cm}{0.805cm}}{\pgfqpoint{11.584cm}{0.795cm}}
\pgfpathcurveto{\pgfqpoint{11.727cm}{0.77cm}}{\pgfqpoint{11.971cm}{0.723cm}}{\pgfqpoint{11.976cm}{0.452cm}}
\pgfpathcurveto{\pgfqpoint{11.983cm}{0.056cm}}{\pgfqpoint{11.539cm}{0.033cm}}{\pgfqpoint{11.358cm}{0.031cm}}
\pgfpathcurveto{\pgfqpoint{11.156cm}{0.028cm}}{\pgfqpoint{10.84cm}{0.057cm}}{\pgfqpoint{10.676cm}{0.366cm}}
\pgfpathlineto{\pgfqpoint{10.956cm}{0.445cm}}
\pgfusepath{fill}

\end{pgfscope}
\end{pgfpicture}
}}; % 13.2 cm width
     \node[anchor=west,xshift=0.745\paperwidth,yshift=0.516\headerheight]{\includegraphics[width=0.171\paperwidth]{figures/Avancez_bg.pdf}}; % 3.6cm
   \end{tikzpicture}};
\end{tikzpicture}

\ \vfill
\makeatletter\@ifundefined{coverfigure}{}{
 \begin{center}
  \includegraphics[width=\textwidth,height=0.4\paperheight,keepaspectratio]{\coverfigure}
 \end{center}
}\makeatother
{\huge\noindent \mytitle\par} % title, 21
{\large\noindent \ \par} % subtitle, 16
%{\normalsize\noindent\textit{Master of Science Thesis}\par};
{\normalsize\noindent\textit{\thesis\ in \nameofprogramme}\par} % 14
\vskip 5mm
{\Large\noindent \uppercase\expandafter{\authors}\par}% 18
\vskip 6mm
{\small\noindent
Department of \departmentname\\
\textit{Division of \divisionnameA}
\makeatletter\@ifundefined{divisionnameB}{}{\textit{and Division of \divisionnameB}}\makeatother\\
\uppercase\expandafter{\university}\\
\city, \country\ \writtenyear\\
\thesis\ \YYYYNN\par}

\newpage
\thispagestyle{empty}
\mbox{}

\newpage
\thispagestyle{empty}
\begin{center}
 {\uppercase\expandafter{\thesis}\ \YYYYNN\par}
 \vskip 40mm
 {\Large \mytitle\par}
 \vskip 5mm
 \makeatletter\@ifundefined{mysubtitle}{}{{\mysubtitle\par}}\makeatother
 \vskip 5mm
 {\thesis\ in \nameofprogramme\par} % (if applicable) (if applicable)
 {\uppercase\expandafter{\authors}\par}
 \vfill
 {Department of Applied Mechanics\par}
 {\textit{Division of \divisionnameA}
  \makeatletter\@ifundefined{divisionnameB}{}{\textit{and Division of \divisionnameB}}\makeatother
 \par}
 {\uppercase\expandafter{\university}\par}
 \vskip 2mm
 {\city, \country\ \writtenyear\par}
\end{center}

\newpage
{\noindent \mytitle\\
\makeatletter\@ifundefined{mysubtitle}{}{\mysubtitle\\}\makeatother
{\uppercase\expandafter{\authors}}\par}
\vskip 10mm
{\noindent \copyright {\uppercase\expandafter{\authorsc}}, \writtenyear\par}
\vskip 20mm
{\noindent \thesis\ \YYYYNN\\
ISSN \ISSN\\
Department of \departmentname\\
Division of \divisionnameA
\makeatletter\@ifundefined{divisionnameB}{}{ and Division of \divisionnameB}\makeatother
\\
\university\\
\postcode\ \city\\
\country\\
Telephone: \telephone\par}
\vfill
\makeatletter\@ifundefined{covercaption}{}{{\noindent Cover:\\\covercaption\par}\vskip 5mm}\makeatother
{\noindent \nameofprinters 
%/ Department of \departmentname
\\\city, \country\ \writtenyear\par}
\thispagestyle{empty}

\newpage
%\thispagestyle{justpage}
\setcounter{page}{1}
\pagenumbering{Roman}
{\noindent \mytitle\\
\makeatletter\@ifundefined{mysubtitle}{}{\mysubtitle\\}\makeatother
\thesis\ in \nameofprogramme\\ % (if applicable)
{\uppercase\expandafter{\authors}}\\
Department of \departmentname\\
Division of \divisionnameA
\makeatletter\@ifundefined{divisionnameB}{}{ and Division of \divisionnameB}\makeatother\\
\university\par}

\phantomsection\addcontentsline{toc}{section}{Abstract}\begin{abstract}
It is important to study the aerodynamic effects on high speed trains, due to both comfort and stability.
The Swedish high speed trains are aiming to go at a speed of 250 km/h. The present work closely connects the aerodynamic effects with the vibration dynamics within the train. 

Two scenarios are simulated, two trains meeting each other and a train leaving a tunnel and is hit by a strong wind gust (35 m/s). From the aerodynamic part, computational fluid dynamics (CFD) is used with the $k$-$\zeta$-$f$ turbulence model. To simulate both scenarios a moving mesh needs to be used. From the CFD the moments and forces from the pressure and traction on the train body is calculated, and these loads are taking into a low order mathematical model that simulates the vibration dynamics in the train.

For the scenario with meeting trains the train experiences some slight vibrations, causing discomfort, but has no impact on the stability of the train, but for the scenario with the tunnel and side wind the train has a very high risk of derailment. 

From the results of the dynamics simulations a comfort and stability measurements were constructed based on the vibrations in the train car and the risk of wheel climbing. From simulating different speeds of the train it could be seen that the comfort and stability changes linearly with the speed.
Work was also done to see how much impact the coupling between the car bodies can have on the comfort and stability. A comparison was made to a simple, stiff coupling and one that optimizes a set of passive dampers. It's seen that the coupling can make a difference of around 8\% in comfort and stability, which equals an effect of lowering the speed 3-5 m/s.

Semi-active dampers of sky-hook and ground-hook type were also tested in the coupling, but they only showed insignificant changes or changes to the worse.
\end{abstract}
\noindent Keywords: \keywords

% Either a swedish translation, or an emtpy page.
\newpage
\mbox{}

\newpage
\phantomsection\addcontentsline{toc}{section}{Contents}
\tableofcontents

% Here you can add preface and notations
\cleardoublepage
\phantomsection\addcontentsline{toc}{section}{Preface}\section*{Preface}
In this study the forces from aerodynamic simulations have been used to analyze the vibrational behaviour of a high speed train in respect to comfort and stability. A scenario of two trains meeting and a scenario of a train leaving a tunnel with a strong side wind are simulated.
The work has been carried out from September 2008 to April 2009 at the Department of \departmentname, Division of \divisionnameA\ and Division of \divisionnameB, Chalmers University of Technology, Sweden, with Mikael \"Ohman as student and Professor Viktor Berbyuk as supervisor for the vibration dynamics and Erik Bjerklund as student and Associate Professor Sini\v{s}a Krajnovi\'{c} as supervisor for the fluid mechanics.

\section*{Aknowledgements}
We would like to say our thanks to our supervisors and the department of \departmentname\ for their support, to Albin Jonsson for his help with the vibration dynamics and to the CFD support at AVL, Dr. Branislav Basara, Andreas Diemath and J\"urgen Schneider. A special thanks to AVL for providing the Fire licenses.
Without their help this project would not have been possible.
\vskip 1cm
\noindent \city\ \writtenmonth\ \writtenyear\\
\authorsc
\newpage
%\phantomsection\addcontentsline{toc}{section}{Notations}\input{Notations}

\cleardoublepage
\setcounter{page}{1}
\pagenumbering{arabic}

% Real contents of report starts here
% Splitting it up to several files help when working together.
% Floatbarriers prevent figures from beeing placed into the next chapter.
\section{Introduction}
% \subsection{Background}
Due to the magnitude of the velocity of a high speed train (HST) the flow around it becomes an increasingly important factor. 
The unsteadiness and the different forces and moments that start acting on the train have a big impact on the stability of the train 
and safety and comfort of the passengers. This can probably be countered using damping in the bogies and coupling between the trains.
This is the reason behind this project. To find the forces and moments affecting the train 
and if it's the effects on discomfort and instability can be reduced.
In previous studies of HST instabilities it has either been purely aerodynamic or dynamic with simplified aerodynamic forces.
In this case it is intended to combine them to get an as accurate simulation as possible. 

\subsection{Purpose}
The purpose of this project is to combine two different fields, 
aerodynamics and dynamics, to accurately simulate the HST by
considering all the forces acting on the train in order to get an accurate model.
The goal of this project is to combine the aerodynamic forces that acts on a HST 
with the dynamics on the trains bogie and coupling. 
An interesting final result would be to find how sensitive the dynamics in the train is to varying velocity.

From vehicle dynamics side the focus of the project will be two folds:
\begin{itemize}
\item Create low-order mathematical and computational models for vibration dynamics and stability analysis of HST taking into account 
aerodynamic excitations. 
The models must implement the conventional bogie and conventional car-body coupling functional component mechanical models.
\item Using created models simulate the vibration dynamics and study stability of motion of HST for the two scenarios.
\end{itemize}

\subsection{Limitations}
Two different scenarios were studied. One where two trains meet each other and one where a single train has left a tunnel and is met by a strong wind gust. The model of the train used for the CFD simulations was an ICE2 train. The train consists of two locomotives and one car in the middle making it symmetrical. It has bogies and inter-car gaps. The parameters for the dynamics model were that of a typical HST. For the meeting trains the simulation was carried out at three different velocities, 67 m/s, 70 m/s and 73 m/s, which correspond to 240, 250 and 260 km/h respectively. For the scenario where the train coming out of a tunnel the speed of the train was 70 m/s and the speed of the wind gust was 35 m/s.

\subsection{Approach}
The aerodynamics was simulated with AVL-Fire CFD solver. For the double train, two train models were used with a moving mesh to simulate the two trains meeting each other. For the second scenario with the HST coming out of the tunnel a model of a ICE2 train in a simulated windtunnel with a moving crosswind was used. The forces and moments acting on the each train car on one of the trains were then handed over for the dynamics calculations.

The dynamics were simulated with a low order mathematical model using MATLAB and functional components for bogie and coupling.
The input were the forces and moments calculated from the CFD simulations and the forces acting on the wheels from the rail.
Unknown parameters, such as damping coefficients in the bogie, were decided by minimizing a cost function for comfort and stability.
As a last step, active damping was added to the model.\FloatBarrier 

\subsection{Theory}
The contact forces are calculated from creep in the contact plane.
The creep is defined as 
\begin{align}
 \xi_x    &= \frac{V_{wheel,x}}{V_0}\\
 \xi_\eta &= \frac{V_{wheel,\eta}}{V_0}
\end{align}
where $x$ is in the longitudinal direction, $\eta$ is perpendicular $x$ and the normal at the contact point and $V_0$ is the forward velocity of the train.

%Is it the purpose of these derivitions to not assume small deflections, as it is already solved as a nonlinear force.
\begin{figure}[htpb!]
 \centering
  \includegraphics[width=0.8\textwidth]{figures/QQ_0409.pdf}
 \caption{A plot!}
 \label{fig:plot} % Labels must come after caption!
\end{figure}

\begin{figure}[htpb!]
\centering
%\begin{tikzpicture}
%  \begin{axis}[xlabel=$x$,ylabel=$y$]
%   \addplot[color=blue,mark=*] plot coordinates { (0,2) (2,3) (3,1) };
%   \addlegendentry{Case 1}
%   \addplot[color=red,mark=x] plot coordinates { (0,0) (1,1) (2,1) (3,2) };
%   \addlegendentry{Case 2}
%  \end{axis}
% \end{tikzpicture}
 \caption{A tikz plot!}
 \label{fig:tikz_plot}
\end{figure}

Right side refers to the right in the forward direction of the train, i.e. the right wheel is not shown in figure \ref{fig:plot}.

The velocities can now be calculated as
% \begin{align}
% \omega_0 &= \frac{V_0}{r_0}\\
%  s       &= \begin{cases}+1&\text{if right side}\\-1&\text{if left side}\end{cases}\\
%  r       &= r[-s x_2]\\
%  w       &= s w_0 - x_2\\
%  \uv{d}  &= \begin{pmatrix}0\\w\\r\end{pmatrix}\\
%  \um{R}_z     &= \begin{pmatrix}
%         \cos[-\varphi_3] & -\sin[-\varphi_3] & 0\\
%         \sin[-\varphi_3] &  \cos[-\varphi_3] & 0\\
%         0                &  0                & 1
%      	\end{pmatrix} \\
%  \uv{V}_{wheel} &= \uv{\dot{x}} +\begin{pmatrix}V_0\\0\\0\end{pmatrix} + \um{R}_z \cdot (\uv{d} \times (\uv{\dot{\varphi}}+\begin{pmatrix}0\\\omega_0\\0\end{pmatrix}))
% \end{align}
where the $r$ is an interpolating function of the geometry of a standardized S1002 wheel profile \cite{railvehicledynamics_eandersson}.

\begin{table}[htpb!]
 \centering
 \caption{Explanation of variables for calculating the creep.}
 \begin{tabular}{c | l}
%  $\uv{x}$       & Translation vector of wheelset.\\
%  $\uv{\varphi}$ & Rotation vector of wheelset.\\
  $r_0$		 & Wheel radius at nominal contact point.\\
  $r$            & Wheel radius at current contact point.\\
  $w_0$		 & Lateral distance from COG to nominal contact point.\\
  $w$		 & Lateral distance from COG to current contact point.\\
  $V_0$		 & Nominal forward velocity of train.\\
  $\omega_0$	 & Nominal angular velocity of wheels.\\
  $\delta$	 & Inclination of wheels (positive).\\
  $s$		 & Side dependence, +1 for right side, -1 for left side.
 \end{tabular}
\end{table}

\subsubsection{Contact Patch Dimensions}
To calculate the forces the dimensions of the contact patch are required.
According to Hertz theory the dimensions are calculated as \cite{hallfasthetslara_bsundstrom} and \cite{railvehicledynamics_eandersson}
\begin{align}
 A+B    =& \frac12 \left(\frac{1}{r_{\eta,r}}+\frac{1}{r_{\eta,w}}+\frac{1}{r_{x,r}}+\frac{1}{r_{x,w}} \right)\\
 \nonumber B-A =& \frac12 \left(\left(\frac{1}{r_{x,r}}-\frac{1}{r_{\eta,r}}\right)^2 + 
                          \left(\frac{1}{r_{x,w}}-\frac{1}{r_{\eta,w}}\right)^2 +\right.\\
               &\left.   2\left(\frac{1}{r_{x,r}}-\frac{1}{r_{\eta,r}}\right)
                          \left(\frac{1}{r_{x,w}}-\frac{1}{r_{\eta,w}}\right)\cos[2\phi_z]\right)^\frac12\\
 \Psi   =& \sqrt[3]{\frac{3 N}{2 (A+B)} \left(\frac{1-\nu_w^2}{E_w} +\frac{1-\nu_r^2}{E_r}\right)}\\ 
 \theta =& \arccos\left[\frac{B-A}{A+B}\right]\\
 a      =& m \Psi\\
 b      =& n \Psi
\end{align}
where $m = m[\theta]$ and $n = n[\theta]$. There is no explicit expression for $m[\theta]$ or $n[\theta]$ but they can be solved numerically with (4.25) through (4.32) in \cite{contact_mechanics_kljohnson}. This has been done in table 8-1 in \cite{railvehicledynamics_eandersson} which values are used.

\begin{table}[htpb!]
 \centering
 \caption{Explanation of variables for calculating the creep.}
 \begin{tabular}{c | l}
  $N$                                     & Normal force in contact.\\
  $E_w,E_r$                               & Modulus of elasticity in the wheel and rail.\\
  $\nu_w,\nu_r$                           & Poissons ratio for the wheel and rail.\\
  $a,b$                                   & Contact patch dimension along $x$ and $\eta$.\\
  $r_{x,w},r_{x,r},r_{\eta,w},r_{\eta,r}$ & Radius around the respective axis for wheel and rail.\\
  $\phi_z$                                & The angle of the wheel.
 \end{tabular}
\end{table}
The radius $r_{\eta,r}$ is zero since the rail is straight. 
The other radii varies with the contact point and wheel profile. However they are approximated as constant for the contact geometry at a point far from the flange where the wheel is almost flat. In that case $r_{x,w}=\infty$ and $r_{\eta,w} = r$.
For the wheel the radius over the top of the rail is almost constant and be approximated as $r_{x,r} = 300$mm.
In this case $\eta$-axis is parallell with the $y$-axis.
The effect from the rotation of the wheel is also neglected, as this angle will always be small, i.e. linearisation of $cos[2\phi_z] \approx 1$. This simplifies the equations and with the approximations of the radii the variables $m$ and $n$ are constant.

The rest of the material parameters were set to $E_w = E_r = 200$ GPa and $\nu_w = \nu_r = 0.3$.

\end{document}
\subsubsection{Wheel Forces from Creep}
The creep forces are calculated with Kalker's linear theory \cite{railvehicledynamics_eandersson}.
The coefficients for the creep, $\uv{C} = \uv{C}[\frac{a}{b}]$, are interpolated from a set of data taken from table 8-2 in \cite{railvehicledynamics_eandersson}.

\begin{align}
 \uv{F}_{lin} &= - G a b \begin{pmatrix}
                          C_{11} \xi_x\\C_{22} \xi_\eta
                         \end{pmatrix}\\
 u            &= \frac{|\uv{F}_{lin}|}{\mu N}\\
 \uv{\hat{F}} &= \uv{F}_{lin} \begin{cases}
                  1-\frac{u}{3}+\frac{u^2}{27} &\text{if $u < 3$} \\
                  \frac{\mu N }{|\uv{F}_{lin}|} & \text{otherwise}
                 \end{cases}
\end{align}

The behaviour of the flange will be greatly simplified as done in \cite{appwheelrailcurvedtracks_jpombo}
\begin{align}
 F_{flange} = \begin{cases}
               k_{flange}( n_{flange}-x_2) & \text{if } x_2 >  n_{flange} \wedge s > 0 \text{ (right side)}\\
               k_{flange}(-n_{flange}-x_2) & \text{if } x_2 < -n_{flange} \wedge s < 0 \text{ (left side)}\\
               0                           & \text{otherwise}
              \end{cases}
\end{align}
where $F_{flange}$ acts in the $y$-axis and is added to  $F_y$.

The force vector $\uv{\hat{F}}$ is here acting on the contact plane, so to get the horizontal component
\begin{align}
 \uv{F} = \begin{pmatrix}
           F_x\\
           F_y
          \end{pmatrix} =
          \begin{pmatrix}
           \hat{F}_1\\
           \hat{F}_2\frac{\delta}{\sqrt{1+\delta^2}}+F_{flange}
          \end{pmatrix}
\end{align}

The contributions to the nonlinear force on the system will be
\begin{align}
 \uv{R}_e = \begin{pmatrix}
      F_x\\F_y\\0\\ 0\\r F_x\\-w F_x -w F_y \sin[\phi_3]
     \end{pmatrix}
\end{align}
which is assembled into $\uv{R}$ for each wheel.
\FloatBarrier
% \input{Method}\FloatBarrier
% \input{Results}\FloatBarrier
% \input{Conclusions}
% \input{Recommendations}

% And the bilbiography saved as mybib.bib
\bibliography{mybib}{}
\bibliographystyle{plain}

% Appendices

\end{document}

